\chapter{Projet de programmation}
\label{chap:programmation}

\section{Pr\'esentation du projet de programmation}

La v\'erification de la s\'ecurit\'e d'un site web constitue un probl\`eme majeur pour
les entreprises, \'ecoles et administrations b\'eninoises, particuli\`erement dans un
contexte o\`u l'expertise en cybers\'ecurit\'e reste rare et co\^uteuse. Les m\'ethodes
disponibles aujourd'hui -- audits professionnels payants, outils techniques
disperss\'es et en anglais, ou absence totale de v\'erification -- ne r\'epondent pas
au besoin d'un propri\'etaire de site qui veut simplement savoir si son site est
s\^ur et quoi corriger. \textbf{CyberGuard B\'enin} est une plateforme web con\c{c}ue
pour rendre ce diagnostic accessible, automatique et compr\'ehensible.

\section{Pr\'esentation de la solution}

Pour r\'epondre \`a la probl\'ematique de d\'emocratisation de l'audit de s\'ecurit\'e web,
nous proposons le d\'eveloppement d'une plateforme compl\`ete utilisant les
technologies Laravel et React. L'utilisateur cr\'ee un compte, saisit l'URL de son
site, et la plateforme lance une s\'erie de v\'erifications non intrusives (HTTPS,
certificat SSL, en-t\^etes de s\'ecurit\'e, exposition d'informations serveur,
cookies). Le syst\`eme traduit ces r\'esultats techniques en un score sur 100 et un
niveau de risque (Bon / Moyen / Faible), accompagn\'e de recommandations concr\`etes
et compr\'ehensibles. Un rapport est g\'en\'er\'e, et l'utilisateur peut suivre
l'\'evolution de la s\'ecurit\'e de son site dans le temps via un tableau de bord et
un historique.

\section{Analyse et mod\'elisation}

\subsection{Choix technique}

La mod\'elisation de ce projet s'appuie sur les diagrammes UML standard (cas
d'utilisation, classes, objets, s\'equence), qui permettent de visualiser les
acteurs, les interactions et la structure statique du syst\`eme avant son
impl\'ementation. Le tableau~\ref{tab:stack} r\'esume les choix technologiques
retenus.

\begin{table}[H]
  \centering
  \caption{Stack technique retenue}
  \label{tab:stack}
  \begin{tabular}{ll}
    \toprule
    \textbf{Couche} & \textbf{Choix} \\
    \midrule
    Frontend & React (Vite) \\
    Backend & Laravel \\
    Base de donn\'ees & MySQL (MariaDB via XAMPP) \\
    Authentification & Laravel Sanctum \\
    G\'en\'eration PDF & DomPDF \\
    \bottomrule
  \end{tabular}
\end{table}

\subsection{\'Etude de l'existant}

Des outils comme \textbf{Mozilla Observatory} et \textbf{SSL Labs} existent d\'ej\`a
pour analyser la s\'ecurit\'e d'un site web. Ils sont reconnus et techniquement
solides, en s'appuyant sur des v\'erifications approfondies des en-t\^etes de
s\'ecurit\'e et des certificats SSL/TLS.

\subsection{Insuffisances}

Ces outils pr\'esentent toutefois plusieurs limites dans le contexte vis\'e par
notre projet. D'abord, ils sont exclusivement en anglais, ce qui constitue un
frein r\'eel pour une partie du public b\'eninois. Ensuite, leurs r\'esultats restent
tr\`es techniques~: un propri\'etaire de site non-technicien obtient une note ou une
liste d'en-t\^etes manquants, sans traduction claire de ce que cela signifie
concr\`etement ni de la marche \`a suivre pour corriger le probl\`eme. Enfin, ces
outils sont dispers\'es~: v\'erifier HTTPS, le certificat SSL et les en-t\^etes de
s\'ecurit\'e n\'ecessite souvent de passer par plusieurs services diff\'erents, sans
suivi centralis\'e ni historique des analyses pass\'ees. CyberGuard B\'enin se
positionne comme une alternative qui centralise ces v\'erifications, restitue un
score unique et des recommandations en fran\c{c}ais compr\'ehensibles par un public
non-technicien, tout en conservant un historique consultable dans le temps.

\subsection{Repr\'esentation de la mod\'elisation}

\subsubsection{Diagramme des cas d'utilisation}

\paragraph{Identification des acteurs}

Trois acteurs interagissent avec le syst\`eme (tableau~\ref{tab:acteurs}).

\begin{table}[H]
  \centering
  \caption{Identification des acteurs}
  \label{tab:acteurs}
  \begin{tabular}{lp{9.5cm}}
    \toprule
    \textbf{Acteur} & \textbf{Description} \\
    \midrule
    Utilisateur & Personne disposant d'un compte sur la plateforme, propri\'etaire ou responsable d'un site web qu'elle souhaite analyser. \\
    Administrateur & Personne charg\'ee de la gestion de la plateforme (comptes, statistiques, param\`etres de scoring). H\'erite de tous les droits de l'Utilisateur. \\
    Syst\`eme & Acteur non-humain~: partie automatis\'ee de l'application qui ex\'ecute les v\'erifications techniques et produit les r\'esultats. \\
    \bottomrule
  \end{tabular}
\end{table}

\paragraph{Identification des cas d'utilisation}

Pour chaque acteur, les cas d'utilisation suivants ont \'et\'e identifi\'es~:

\textbf{Utilisateur}
\begin{itemize}[nosep]
  \item G\'erer son compte (inscription, modification, r\'einitialisation du mot de passe)
  \item S'authentifier (connexion par email/mot de passe, ou via un compte Google d\'ej\`a existant -- Google n'est jamais un moyen d'inscription)
  \item Lancer une analyse
  \item Consulter les rapports (historique, d\'etail, t\'el\'echargement PDF)
  \item Consulter le tableau de bord
  \item \^Etre notifi\'e de la fin d'une analyse (notification in-app lorsqu'une analyse lanc\'ee en arri\`ere-plan se termine)
\end{itemize}
\textbf{Administrateur}
\begin{itemize}[nosep]
  \item G\'erer les comptes utilisateurs (activer/d\'esactiver, promouvoir/r\'etrograder le r\^ole, supprimer)
  \item Mod\'erer les analyses (consulter et supprimer l'analyse de n'importe quel utilisateur)
  \item Administrer la plateforme (statistiques globales, param\`etres de scoring)
\end{itemize}
\textbf{Syst\`eme}
\begin{itemize}[nosep]
  \item Traiter l'analyse (v\'erifications techniques, calcul du score, g\'en\'eration du rapport, notification de l'utilisateur)
\end{itemize}

\paragraph{Tableau r\'ecapitulatif des cas d'utilisation}

\begin{longtable}{clc}
  \caption{Tableau r\'ecapitulatif des cas d'utilisation}
  \label{tab:recap-cu} \\
  \toprule
  \textbf{\#} & \textbf{Cas d'utilisation} & \textbf{Acteur} \\
  \midrule
  \endfirsthead
  \toprule
  \textbf{\#} & \textbf{Cas d'utilisation} & \textbf{Acteur} \\
  \midrule
  \endhead
  1 & G\'erer son compte & Utilisateur \\
  2 & S'authentifier & Utilisateur \\
  3 & Lancer une analyse & Utilisateur \\
  4 & Consulter les rapports & Utilisateur \\
  5 & Consulter le tableau de bord & Utilisateur \\
  6 & \^Etre notifi\'e de la fin d'une analyse & Utilisateur \\
  7 & G\'erer les comptes utilisateurs & Administrateur \\
  8 & Mod\'erer les analyses & Administrateur \\
  9 & Administrer la plateforme & Administrateur \\
  10 & Traiter l'analyse & Syst\`eme \\
  \bottomrule
\end{longtable}

\paragraph{Diagramme proprement dit}

Le diagramme ci-dessous formalise les acteurs et cas d'utilisation identifi\'es,
avec une relation de g\'en\'eralisation (Administrateur h\'erite d'Utilisateur) et
une relation d'inclusion (\og~Lancer une analyse~\fg{} inclut \og~Traiter
l'analyse~\fg).

\TODOFigure{diagramme-cas-utilisation.jpeg}{Diagramme de cas d'utilisation}{fig:cas-utilisation}

\paragraph{Description textuelle de cinq cas d'utilisation}

\subparagraph{Cas d'utilisation~: G\'erer les comptes utilisateurs}
\begin{table}[H]
  \centering
  \small
  \begin{tabular}{p{3.2cm}p{9.5cm}}
    \toprule
    Acteur principal & Administrateur \\
    Objectif & Permettre \`a un administrateur de g\'erer les comptes des utilisateurs de la plateforme (statut, r\^ole). \\
    Pr\'econdition & L'administrateur est authentifi\'e. \\
    Sc\'enario nominal &
      1. Acc\`es \`a la liste des comptes utilisateurs.
      2. Recherche par nom/email, filtre par statut.
      3. Activation ou d\'esactivation d'un compte.
      4. Promotion ou r\'etrogradation du r\^ole (utilisateur $\leftrightarrow$ administrateur).
      5. Suppression d'un compte si n\'ecessaire. \\
    Postcondition & L'\'etat du compte cibl\'e (statut, r\^ole, ou existence) est mis \`a jour en cons\'equence. \\
    Exception & Tentative d'un administrateur de d\'esactiver, r\'etrograder ou supprimer son propre compte $\rightarrow$ refus\'ee (garde-fou anti-auto-modification). Tentative d'acc\`es par un utilisateur non administrateur $\rightarrow$ acc\`es refus\'e. \\
    \bottomrule
  \end{tabular}
\end{table}

\subparagraph{Cas d'utilisation~: Lancer une analyse}
\begin{table}[H]
  \centering
  \small
  \begin{tabular}{p{3.2cm}p{9.5cm}}
    \toprule
    Acteur principal & Utilisateur \\
    Objectif & Permettre \`a l'utilisateur de demander l'analyse de s\'ecurit\'e d'un site web. \\
    Pr\'econdition & L'utilisateur est authentifi\'e. \\
    Sc\'enario nominal &
      1. L'utilisateur saisit l'URL du site \`a analyser.
      2. Le syst\`eme valide le format de l'URL.
      3. Le syst\`eme d\'eclenche le cas d'utilisation \og~Traiter l'analyse~\fg.
      4. Une confirmation de lancement est affich\'ee \`a l'utilisateur. \\
    Postcondition & Une analyse est en cours de traitement pour le site indiqu\'e. \\
    Exception & URL invalide ou site inaccessible $\rightarrow$ message d'erreur affich\'e, aucune analyse n'est lanc\'ee. \\
    \bottomrule
  \end{tabular}
\end{table}

\subparagraph{Cas d'utilisation~: Traiter l'analyse}
\begin{table}[H]
  \centering
  \small
  \begin{tabular}{p{3.2cm}p{9.5cm}}
    \toprule
    Acteur principal & Syst\`eme \\
    Objectif & Ex\'ecuter automatiquement les v\'erifications de s\'ecurit\'e et produire un r\'esultat exploitable. \\
    Pr\'econdition & Le cas \og~Lancer une analyse~\fg{} a \'et\'e d\'eclench\'e avec une URL valide. \\
    Sc\'enario nominal &
      1. V\'erification que l'h\^ote ne pointe pas vers une ressource r\'eseau priv\'ee ou interne.
      2. V\'erification de la pr\'esence du HTTPS.
      3. Contr\^ole de la validit\'e du certificat SSL.
      4. Analyse des en-t\^etes de s\'ecurit\'e HTTP.
      5. V\'erification de la configuration des cookies (Secure, HttpOnly, SameSite).
      6. Recherche de vuln\'erabilit\'es simples et non intrusives.
      7. Recherche de fichiers sensibles expos\'es publiquement (.env, .git...).
      8. V\'erification de la protection du domaine contre l'usurpation d'email (SPF/DMARC).
      9. Calcul du score sur 100 et du niveau de risque.
      10. G\'en\'eration des recommandations, en langage compr\'ehensible par un non-technicien.
      11. G\'en\'eration du rapport PDF.
      12. Enregistrement du r\'esultat en base de donn\'ees.
      13. Notification de l'utilisateur si l'analyse \'etait suivie en arri\`ere-plan (\og~\^Etre notifi\'e de la fin d'une analyse~\fg). \\
    Postcondition & Un rapport d'analyse (score, risque, recommandations, PDF) est disponible et associ\'e \`a l'utilisateur. \\
    Exception & Adresse non autoris\'ee, site inaccessible ou timeout $\rightarrow$ l'analyse est marqu\'ee en \'echec avec un motif explicite, l'utilisateur est notifi\'e. \\
    \bottomrule
  \end{tabular}
\end{table}

\subparagraph{Cas d'utilisation~: Consulter les rapports}
\begin{table}[H]
  \centering
  \small
  \begin{tabular}{p{3.2cm}p{9.5cm}}
    \toprule
    Acteur principal & Utilisateur \\
    Objectif & Permettre \`a l'utilisateur de revoir ses analyses pass\'ees et d'en r\'ecup\'erer le d\'etail. \\
    Pr\'econdition & L'utilisateur est authentifi\'e et poss\`ede au moins une analyse termin\'ee. \\
    Sc\'enario nominal &
      1. L'utilisateur acc\`ede \`a son historique d'analyses.
      2. Le syst\`eme affiche la liste des analyses (date, site, score, niveau de risque).
      3. L'utilisateur s\'electionne une analyse.
      4. Le syst\`eme affiche le d\'etail du rapport.
      5. L'utilisateur peut t\'el\'echarger le rapport au format PDF. \\
    Postcondition & L'utilisateur a consult\'e ou t\'el\'echarg\'e le rapport souhait\'e. \\
    Exception & Aucune analyse enregistr\'ee $\rightarrow$ message indiquant qu'aucun historique n'est disponible. \\
    \bottomrule
  \end{tabular}
\end{table}

\subparagraph{Cas d'utilisation~: Mod\'erer les analyses}
\begin{table}[H]
  \centering
  \small
  \begin{tabular}{p{3.2cm}p{9.5cm}}
    \toprule
    Acteur principal & Administrateur \\
    Objectif & Permettre \`a un administrateur de superviser l'ensemble des analyses de la plateforme, tous utilisateurs confondus, et d'en supprimer une si n\'ecessaire. \\
    Pr\'econdition & L'administrateur est authentifi\'e. \\
    Sc\'enario nominal &
      1. Acc\`es \`a la liste globale des analyses.
      2. Affichage de toutes les analyses avec leur propri\'etaire, ind\'ependamment de l'utilisateur connect\'e.
      3. Recherche par URL ou par propri\'etaire, filtre par niveau de risque.
      4. Consultation du d\'etail d'une analyse.
      5. Suppression de l'analyse (et de son rapport PDF associ\'e) si n\'ecessaire. \\
    Postcondition & La liste refl\`ete l'\'etat courant des analyses~; une analyse supprim\'ee n'est plus consultable. \\
    Exception & Tentative d'acc\`es par un utilisateur non administrateur $\rightarrow$ acc\`es refus\'e. \\
    \bottomrule
  \end{tabular}
\end{table}

\subsubsection{Diagramme des classes}

Un diagramme de classes est un diagramme UML qui repr\'esente la structure
statique d'un syst\`eme logiciel en mod\'elisant les classes avec leurs attributs,
m\'ethodes et les relations qui les lient. Chaque classe est visualis\'ee sous forme
d'un rectangle divis\'e en trois sections (nom, attributs, m\'ethodes). Les relations
entre classes (association, h\'eritage, agr\'egation, d\'ependance) sont
repr\'esent\'ees par diff\'erents types de fl\`eches et de lignes.

\paragraph{Identification des classes}
\begin{itemize}
  \item Utilisateur
  \item Administrateur
  \item Analyse
  \item Controle
  \item ParametreScore
\end{itemize}

\paragraph{R\`egles de gestion}
\textbf{Utilisateur}
\begin{itemize}[nosep]
  \item Un Utilisateur peut lancer plusieurs Analyses.
  \item Un Administrateur h\'erite de tous les droits d'un Utilisateur (g\'en\'eralisation).
\end{itemize}
\textbf{Analyse}
\begin{itemize}[nosep]
  \item Une Analyse est lanc\'ee par un seul Utilisateur.
  \item Une Analyse contient plusieurs Controles (r\'esultats de v\'erification).
  \item Une Analyse a un statut \'evolutif (en\_cours, terminee, echec).
  \item Un Administrateur peut consulter et supprimer n'importe quelle Analyse, sans restriction de propri\'etaire (mod\'eration).
\end{itemize}
\textbf{Controle}
\begin{itemize}[nosep]
  \item Un Controle est toujours li\'e \`a une seule Analyse.
  \item Un Controle a un statut (conforme / non\_conforme) et peut porter une recommandation.
\end{itemize}
\textbf{ParametreScore}
\begin{itemize}[nosep]
  \item Il n'existe qu'une seule instance de ParametreScore (configuration globale, singleton).
  \item Un Administrateur peut configurer les param\`etres de score (\og~configure~\fg).
\end{itemize}

\paragraph{Dictionnaire des donn\'ees}

Le dictionnaire des donn\'ees ci-dessous recense, pour chaque table de la base de
donn\'ees, l'ensemble des champs, leur type SQL, leur taille et leur r\^ole (cf.
\texttt{schema.sql}).

\begin{table}[H]
  \centering
  \caption{Dictionnaire des donn\'ees -- table \texttt{utilisateurs}}
  \begin{tabular}{lllp{6cm}}
    \toprule
    \textbf{Champ} & \textbf{Type} & \textbf{Taille} & \textbf{Description} \\
    \midrule
    id & BIGINT UNSIGNED & -- & Identifiant unique de l'utilisateur \\
    nom & VARCHAR & 150 & Nom complet de l'utilisateur \\
    email & VARCHAR & 190 & Adresse email (identifiant de connexion) \\
    mot\_de\_passe & VARCHAR & 255 & Mot de passe hach\'e \\
    role & ENUM & -- & \texttt{utilisateur} / \texttt{administrateur} \\
    statut & ENUM & -- & \texttt{actif} / \texttt{desactive} \\
    date\_inscription & TIMESTAMP & -- & Date de cr\'eation du compte \\
    \bottomrule
  \end{tabular}
\end{table}

\begin{table}[H]
  \centering
  \caption{Dictionnaire des donn\'ees -- table \texttt{analyses}}
  \begin{tabular}{lllp{6cm}}
    \toprule
    \textbf{Champ} & \textbf{Type} & \textbf{Taille} & \textbf{Description} \\
    \midrule
    id & BIGINT UNSIGNED & -- & Identifiant unique de l'analyse \\
    utilisateur\_id & BIGINT UNSIGNED & -- & R\'ef\'erence vers l'utilisateur propri\'etaire \\
    url & VARCHAR & 255 & URL du site analys\'e \\
    date\_analyse & TIMESTAMP & -- & Date de lancement de l'analyse \\
    score & TINYINT UNSIGNED & -- & Score obtenu sur 100 \\
    niveau\_risque & ENUM & -- & \texttt{bon} / \texttt{moyen} / \texttt{faible} \\
    statut & ENUM & -- & \texttt{en\_cours} / \texttt{terminee} / \texttt{echec} \\
    chemin\_rapport\_pdf & VARCHAR & 255 & Chemin du rapport PDF g\'en\'er\'e \\
    \bottomrule
  \end{tabular}
\end{table}

\begin{table}[H]
  \centering
  \caption{Dictionnaire des donn\'ees -- table \texttt{controles}}
  \begin{tabular}{lllp{6cm}}
    \toprule
    \textbf{Champ} & \textbf{Type} & \textbf{Taille} & \textbf{Description} \\
    \midrule
    id & BIGINT UNSIGNED & -- & Identifiant unique du contr\^ole \\
    analyse\_id & BIGINT UNSIGNED & -- & R\'ef\'erence vers l'analyse concern\'ee \\
    type & VARCHAR & 50 & https, certificat\_ssl, en\_tetes\_securite, cookies, vulnerabilites, fichiers\_exposes, email\_securise \\
    statut & ENUM & -- & \texttt{conforme} / \texttt{non\_conforme} \\
    description & TEXT & -- & D\'etail du constat \\
    recommandation & TEXT & -- & Recommandation de correction (si non conforme) \\
    \bottomrule
  \end{tabular}
\end{table}

\begin{table}[H]
  \centering
  \caption{Dictionnaire des donn\'ees -- table \texttt{parametres\_score}}
  \begin{tabular}{lllp{6cm}}
    \toprule
    \textbf{Champ} & \textbf{Type} & \textbf{Taille} & \textbf{Description} \\
    \midrule
    id & TINYINT UNSIGNED & -- & Identifiant fixe (toujours 1, ligne unique) \\
    poids\_https & TINYINT UNSIGNED & -- & Pond\'eration du crit\`ere HTTPS \\
    poids\_certificat\_ssl & TINYINT UNSIGNED & -- & Pond\'eration du crit\`ere certificat SSL \\
    poids\_headers & TINYINT UNSIGNED & -- & Pond\'eration du crit\`ere en-t\^etes \\
    poids\_cookies & TINYINT UNSIGNED & -- & Pond\'eration du crit\`ere cookies \\
    poids\_vulnerabilites & TINYINT UNSIGNED & -- & Pond\'eration du crit\`ere vuln\'erabilit\'es \\
    poids\_fichiers\_exposes & TINYINT UNSIGNED & -- & Pond\'eration du crit\`ere fichiers expos\'es \\
    poids\_email\_securise & TINYINT UNSIGNED & -- & Pond\'eration du crit\`ere s\'ecurit\'e email (SPF/DMARC) \\
    seuil\_bon & TINYINT UNSIGNED & -- & Seuil de score pour le niveau \og~bon~\fg \\
    seuil\_moyen & TINYINT UNSIGNED & -- & Seuil de score pour le niveau \og~moyen~\fg \\
    \bottomrule
  \end{tabular}
\end{table}

\paragraph{Diagramme de classes}

\TODOFigure{diagramme-classes.jpeg}{Diagramme de classes}{fig:classes}

\subsubsection{Diagramme d'objet}

Le diagramme d'objets ci-dessous illustre une instanciation concr\`ete du
diagramme de classes~: un utilisateur ayant lanc\'e une analyse termin\'ee avec
trois contr\^oles, et le bar\`eme de score en vigueur, configur\'e par un
administrateur.

\TODOFigure{diagramme-objets.jpeg}{Diagramme d'objets}{fig:objets}

\subsubsection{Diagrammes de s\'equence}

Un diagramme de s\'equence par cas d'utilisation d\'ecrit ci-dessus, entre
l'Utilisateur, le Syst\`eme et la base de donn\'ees.

\paragraph{G\'erer les comptes utilisateurs}

Liste des comptes avec recherche et filtre par statut, puis modification
optionnelle (statut, r\^ole ou suppression), prot\'eg\'ee par un garde-fou
emp\^echant un administrateur de modifier son propre compte.

\TODOFigure{diagramme-sequence-gerer-comptes.jpeg}{Diagramme de s\'equence -- G\'erer les comptes utilisateurs}{fig:sequence-gerer-comptes}

\paragraph{Lancer une analyse}

Saisie de l'URL, validation du format, cr\'eation de l'analyse (statut
\texttt{en\_cours}) en base et confirmation \`a l'utilisateur, avec le chemin
d'exception \og~URL invalide~\fg. Ce cas d\'eclenche automatiquement \og~Traiter
l'analyse~\fg{} (relation \og~include~\fg).

\TODOFigure{diagramme-sequence-lancer-analyse.jpeg}{Diagramme de s\'equence -- Lancer une analyse}{fig:sequence-lancer}

\paragraph{Traiter l'analyse}

Ex\'ecution des sept contr\^oles (HTTPS, certificat SSL, en-t\^etes de
s\'ecurit\'e, cookies, vuln\'erabilit\'es, fichiers expos\'es, s\'ecurit\'e email
SPF/DMARC), calcul du score et du niveau de risque, g\'en\'eration des
recommandations et du rapport, puis enregistrement en base, avec le chemin
d'exception \og~adresse non autoris\'ee, site inaccessible ou timeout~\fg. Le
traitement est d\'el\'egu\'e \`a une file d'attente~: une fois termin\'e,
l'utilisateur est notifi\'e dans l'application s'il suivait cette analyse
(\og~\^Etre notifi\'e de la fin d'une analyse~\fg).

\TODOFigure{diagramme-sequence-traiter-analyse.jpeg}{Diagramme de s\'equence -- Traiter l'analyse}{fig:sequence-traiter}

\paragraph{Consulter les rapports}

Acc\`es \`a l'historique des analyses, s\'election d'une analyse pour en consulter
le d\'etail, et t\'el\'echargement optionnel du rapport PDF, avec le chemin
d'exception \og~aucune analyse enregistr\'ee~\fg.

\paragraph{Mod\'erer les analyses}

Vue globale r\'eserv\'ee \`a l'Administrateur~: m\^emes principes de recherche et
de filtrage que \og~Consulter les rapports~\fg, mais sans restriction au
propri\'etaire, avec une action de suppression suppl\'ementaire. Acc\`es prot\'eg\'e
par le r\^ole administrateur~: une tentative par un utilisateur normal est
refus\'ee (403).

\section{Outils utilis\'es}

\subsection{Mat\'eriels utilis\'es}

Le d\'eveloppement a \'et\'e r\'ealis\'e sur un ordinateur portable \textbf{HP ProBook
450 G7}, dot\'e d'un processeur \textbf{Intel Core i7-10510U} cadenc\'e \`a
1,80~GHz (jusqu'\`a 2,30~GHz en turbo), de \textbf{16~Go de m\'emoire RAM}, d'un
disque de stockage de 238~Go et d'une carte graphique int\'egr\'ee Intel UHD
Graphics, sous \textbf{Windows 11 Pro}.

\subsection{Frontend}
\begin{itemize}
  \item \textbf{React (Vite)}~: biblioth\`eque JavaScript utilis\'ee pour construire
    une interface utilisateur dynamique en Single Page Application, avec
    \texttt{react-router-dom} pour la navigation et \texttt{axios} pour la
    communication avec l'API.
\end{itemize}

\subsection{Backend (API et logique m\'etier)}
\begin{itemize}
  \item \textbf{Laravel}~: framework PHP servant de base \`a l'API RESTful,
    responsable de la logique m\'etier (moteur de v\'erification de s\'ecurit\'e), de
    l'interaction avec la base de donn\'ees via l'ORM Eloquent, et de la gestion
    de l'authentification via \textbf{Laravel Sanctum}.
\end{itemize}

\subsection{Environnement de d\'eveloppement local}
\begin{itemize}
  \item \textbf{XAMPP}~: environnement logiciel int\'egr\'e fournissant Apache, PHP
    et MySQL (MariaDB) pour le d\'eveloppement local de l'application Laravel.
\end{itemize}

\subsection{Conception d'interface}
\begin{itemize}
  \item \textbf{Stitch} (Google)~: outil de g\'en\'eration de maquettes haute-fid\'elit\'e \`a
    partir de prompts textuels, utilis\'e pour produire les \'ecrans pr\'esent\'es en
    section~\ref{sec:maquettes} \`a partir de l'identit\'e visuelle d\'efinie pour le
    projet (palette de couleurs, typographie, style des composants).
\end{itemize}

\subsection{Outils de test}

Deux niveaux de test compl\'ementaires ont \'et\'e utilis\'es tout au long du
d\'eveloppement.

\paragraph{Tests automatis\'es (backend)} L'API est couverte par une suite de
tests \textbf{PHPUnit}, ex\'ecut\'ee via la commande \texttt{php artisan test}.
Chaque test s'appuie sur le trait \texttt{RefreshDatabase} pour repartir d'une
base de donn\'ees propre, et la biblioth\`eque \textbf{Mockery} est utilis\'ee
pour simuler les r\'eponses du service tiers Google (fa\c{c}ade Socialite) sans
d\'ependre d'une v\'eritable connexion \`a Google lors des tests. Ces tests
couvrent notamment~: l'authentification (classique et via Google), les
autorisations d'acc\`es (propri\'etaire d'une ressource, r\^ole administrateur,
garde-fous anti-auto-modification), la validation des URL soumises \`a
analyse (rejet des adresses r\'eseau priv\'ees ou invalides), et les r\`egles de
gestion des param\`etres de scoring.

\paragraph{Tests manuels de bout en bout (frontend)} L'interface a \'et\'e
v\'erifi\'ee de mani\`ere r\'ep\'et\'ee \`a l'aide de \textbf{Playwright}, un outil
d'automatisation de navigateur~: simulation d'un parcours utilisateur complet
(connexion, lancement d'une analyse, consultation du r\'esultat), v\'erification
de l'absence de d\'ebordement horizontal sur mobile \`a plusieurs largeurs
d'\'ecran, et contr\^ole du bon fonctionnement des fonctionnalit\'es
d'administration. Cette approche a permis de d\'etecter plusieurs anomalies
r\'eelles avant qu'elles n'affectent l'utilisateur final (par exemple, une
session non d\'emarr\'ee sur le chemin de retour de l'authentification Google, ou
une condition de course dans le suivi des notifications en arri\`ere-plan).

\subsection{D\'eploiement (mise en production)}

\`A ce stade du projet, la plateforme n'est pas encore d\'eploy\'ee en production~:
le d\'eveloppement et les tests se font en environnement local via XAMPP
(Apache, MySQL) pour le backend et le serveur de d\'eveloppement Vite pour le
frontend. Le choix d'un h\'ebergeur et la mise en production font partie des
perspectives d'\'evolution du projet (voir conclusion).

\section{Maquettes}
\label{sec:maquettes}

Sur la base de l'identit\'e visuelle d\'efinie pour CyberGuard B\'enin --
\og~Le Sceau~\fg, une identit\'e de confiance institutionnelle~: papier chaud
(\#F5F1E6), bleu marine (\#16304A), accent or (\#B8863A), typographie
Georgia pour les titres, composants aux coins nets et bordures fines plut\^ot
que des ombres port\'ees -- des maquettes haute-fid\'elit\'e ont \'et\'e
r\'ealis\'ees avec \textbf{Stitch} pour chacun des \'ecrans identifi\'es \`a
partir des cas d'utilisation, avant l'impl\'ementation compl\`ete du frontend.
L'identit\'e visuelle a \'et\'e enti\`erement red\'efinie en cours de projet~; les
maquettes ci-dessous refl\`etent la version actuelle.

\TODOFigure{maquette-accueil.png}{Accueil, avec la modale Connexion/Inscription}{fig:maq-accueil}
\TODOFigure{maquette-mot-de-passe-oublie.png}{Mot de passe oubli\'e}{fig:maq-mdp-oublie}
\TODOFigure{maquette-mon-profil.png}{Mon profil}{fig:maq-profil}
\TODOFigure{maquette-tableau-de-bord-mobile.png}{Tableau de bord (variante mobile)}{fig:maq-dashboard-mobile}
\TODOFigure{maquette-tableau-de-bord-desktop.png}{Tableau de bord (variante desktop)}{fig:maq-dashboard-desktop}
\TODOFigure{maquette-nouvelle-analyse.png}{Nouvelle analyse (analyse en cours)}{fig:maq-nouvelle-analyse}
\TODOFigure{maquette-resultat-analyse.png}{R\'esultat d'analyse}{fig:maq-resultat}
\TODOFigure{maquette-historique.png}{Historique des analyses}{fig:maq-historique}
\TODOFigure{maquette-detail-rapport.png}{D\'etail d'un rapport, avec t\'el\'echargement PDF}{fig:maq-detail-rapport}
\TODOFigure{maquette-admin-utilisateurs.png}{Administration -- gestion des utilisateurs}{fig:maq-admin-utilisateurs}
\TODOFigure{maquette-admin-analyses.png}{Administration -- mod\'eration des analyses}{fig:maq-admin-analyses}
\TODOFigure{maquette-admin-parametres.png}{Administration -- statistiques et param\`etres}{fig:maq-admin-parametres}

\begin{table}[H]
  \centering
  \caption{Correspondance maquette / cas d'utilisation}
  \begin{tabular}{p{6cm}p{6.5cm}}
    \toprule
    \textbf{Maquette} & \textbf{Cas d'utilisation associ\'e} \\
    \midrule
    Accueil & S'authentifier (email/mot de passe ou Google), G\'erer son compte (inscription) \\
    Mot de passe oubli\'e & G\'erer son compte (r\'einitialisation) \\
    Mon profil & G\'erer son compte (modification) \\
    Tableau de bord & Consulter le tableau de bord \\
    Nouvelle analyse & Lancer une analyse \\
    R\'esultat d'analyse & Traiter l'analyse (r\'esultat), \^Etre notifi\'e de la fin d'une analyse \\
    Historique & Consulter les rapports (liste) \\
    D\'etail d'un rapport & Consulter les rapports (d\'etail + PDF) \\
    Administration -- utilisateurs & G\'erer les comptes utilisateurs \\
    Administration -- analyses & Mod\'erer les analyses \\
    Administration -- param\`etres & Administrer la plateforme \\
    \bottomrule
  \end{tabular}
\end{table}
